\documentclass[10pt]{article}

\usepackage{amsthm}
\usepackage[all]{xy}
\usepackage{mathrsfs}
\usepackage{amssymb}
\usepackage{enumitem}
\usepackage[usenames,dvipsnames]{xcolor}
\usepackage{xcolor}
\usepackage{tikz-cd}
\usepackage[colorlinks=true]{hyperref}
\usetikzlibrary{patterns}
\usetikzlibrary{arrows.meta,angles,quotes,calc}
\usepackage{comment}
\usepackage{extarrows}
\usepackage{mathtools}
\usepackage{seqsplit}
\usepackage{xstring}
\usepackage[normalem]{ulem}
\usepackage{dutchcal}
\usepackage[margin=1in]{geometry}

\newcommand{\derlog}{{\rm DerLog}}
\newcommand{\logder}{\derlog}

\newcommand{\action}{\:\rotatebox[origin=c]{-90}{$\circlearrowright$}\:}

\newcommand{\heart}{\ensuremath\heartsuit}

\usepackage{pgfplots}


\setlength{\marginparwidth}{2cm}



\setcounter{secnumdepth}{4}
\renewcommand{\theparagraph}{\thesubsection.\arabic{paragraph}}

\numberwithin{equation}{section}



\newcommand{\img}{{\rm Im}}
\newcommand{\valcat}[1]{{\rm Val}(#1)}

\newtheorem{theorem}[paragraph]{Theorem}
\newtheorem{proposition}[paragraph]{Proposition}
\newtheorem{question}[paragraph]{Question}
\newtheorem{conjecture}[paragraph]{Conjecture}
\newtheorem{lemma}[paragraph]{Lemma}
\newtheorem{corollary}[paragraph]{Corollary}
\newtheorem{assumption}[paragraph]{Assumption}
\newtheorem{reduction}[paragraph]{Reduction}

\theoremstyle{definition} 
\newtheorem{definition}[paragraph]{Definition}
\newtheorem{example}[paragraph]{Example}
\newtheorem{warning}[paragraph]{Warning}
\newtheorem{situation}[paragraph]{Situation}
\newtheorem{construction}[paragraph]{Construction}

\theoremstyle{remark}
\newtheorem{remark}[paragraph]{Remark}


\def\id{\mathrm{id}}
\def\Mon{\mathbf{Mon}}
\def\Ext{\operatorname{Ext}}
\def\Hom{\operatorname{Hom}}
\def\Cone{\operatorname{Cone}}
\def\cone{\Cone}
\def\Spec{\operatorname{Spec}}
\def\spec{\Spec}
\def\SSpec{\underline{\operatorname{Spec}}}
\def\Cones{\mathbf{Cones}}
\def\Sets{\mathbf{Sets}}
\def\LogSch{\mathbf{LogSch}}
\def\logGm{\mathbf{G}_{\log}}
\def\Gm{\mathbf{G}_m}
\def\Pic{\mathrm{Pic}}
\def\LogPic{\mathrm{LogPic}}
\def\TroPic{\mathrm{TroPic}}
\DeclareMathOperator*{\colim}{colim}
\def\tropGm{\mathbf{G}_{\rm trop}}
\def\et{\operatorname{\acute{e}t}}
\def\Et{\operatorname{\acute{E}t}}
\def\bExt{\operatorname{\mathbf{Ext}}}
\def\Rh{\operatorname{Rhiz}}
\def\rh{\operatorname{rhiz}}
\def\Opens{\operatorname{Opens}}

\newcommand{\sharpish}[1]{#1^{\# \rm ish}}


\newcommand{\pt}{{\rm pt}}


\let \varinjlim \relax 
\DeclareMathOperator*{\varinjlim}{colim}


\renewcommand{\frak}[1]{\ensuremath{\mathfrak{#1}}}



\newcommand{\norm}[1]{\left\|{#1}\right\|}




\ExplSyntaxOn
\cs_new_eq:NN \Repeat \prg_replicate:nn
\ExplSyntaxOff

\newlist{variantenum}{enumerate}{1}
\setlist[variantenum,1]{label=(\roman*), ref=(\roman*)}

\newcounter{primecnt}[variantenumi]






\newcommand{\face}[1]{{#1}_{\rm face}}
\newcommand{\tropet}[1]{{#1}_{{\rm trop}\et}}



%to control the style of bold text
\newcommand{\mbold}[1]{\mathbf{#1}}
\newcommand{\GG}{\mbold{G}}
\newcommand{\kah}[1]{\Omega_{#1}}
\newcommand{\lkah}[1]{\Omega^\ell_{#1}}
\newcommand{\lccx}[1]{\mbold{L}^\ell_{#1}}
\newcommand{\glccx}[1]{\mbold{L}^G_{#1}}
\def\Glog{\mbold{G}_{\log}}
\def\Gtrop{\mbold{G}_{\rm trop}}
\newcommand{\gp}[1]{#1^{\rm gp}}
\newcommand{\Aff}{\mbold{A}}
\newcommand{\NN}{\mbold{N}}
\newcommand{\OO}{\mathcal{O}}
\newcommand{\usch}[1]{#1^\circ}
\newcommand{\ZZ}{\mbold{Z}}
\newcommand{\CC}{\mbold{C}}
\newcommand{\QQ}{\mbold{Q}}
\newcommand{\RR}{\mbold{R}}
\newcommand{\PP}{\mbold{P}}
\newcommand{\Blog}{Bl}
\newcommand{\HHom}{\underline{\rm Hom}}

\newcommand{\pb}{\ar[dr, phantom, very near start, "\ulcorner"]}
\newcommand{\lpb}{\ar[dr, phantom, very near start, "\ulcorner \ell"]}
\newcommand{\msout}[1]{\text{\sout{\ensuremath{#1}}}}
\newcommand{\lpbstrict}{\ar[dr, phantom, very near start, "\ulcorner \msout{\ell}"]}

\renewcommand{\tilde}[1]{\ensuremath{\widetilde{#1}}}
\renewcommand{\hat}[1]{\ensuremath{\widehat{#1}}}
\newcommand{\scr}[1]{\ensuremath{\mathscr{#1}}}
\newcommand{\af}[1]{\Theta_{#1}}


\newcommand{\utopcone}[1]{\left\vert\Cone #1 \right\vert}
\newcommand{\utopocone}[1]{\left\vert\Cone^\circ #1\right\vert}

\newcommand{\num}[1]{\langle #1 \rangle}
\newcommand{\bra}[1]{\left[ {#1} \right]}
\newcommand{\pra}[1]{\left( {#1} \right)}
\newcommand{\longequals}{\xlongequal{\: \:}}
\newcommand{\longsimeq}{\xrightarrow\sim}

\newcommand{\sharpening}{makeover}

\newcommand{\baroplus}{\overnorm{\oplus}}
\newcommand{\barbigoplus}{\overnorm{\bigoplus}}


\renewcommand{\bar}[1]{\overnorm{#1}\vphantom{#1}}

\renewcommand{\iff}{\Leftrightarrow}

\newcommand{\shp}[1]{#1_\sharp}
\newcommand{\gs}[1]{\gp{\shp{#1}}}
\newcommand{\inr}[1]{#1^\square}
\newcommand{\inrp}[1]{#1^{\square+}}


\newcommand{\smet}[1]{#1_{\acute{e}t}}

\newcommand{\PSh}[1]{\hat{#1}}


\newcommand{\val}[1]{#1^{\rm val}}



%declares a command \cursive that works in text mode that gives a script r
\usepackage{frcursive}
\newcommand{\Tr}[1]{T^{\text{\textit{\cursive{r}}}}_{#1}}
\newcommand{\Nr}[1]{N^{\text{\textit{\cursive{r}}}}_{#1}}

\newcommand{\varv}{v}






\def\overnorm#1{\overline{#1}\vphantom{#1}}

\def\ologGm{{\mathbf G}_{\rm trop}}


\newcommand{\loc}[2]{#1_{(#2)}}

\newcommand{\Log}{{\cal L}og}

\frenchspacing




\begin{document}
\section{Problem statement}

In the following, a monoid is an integral, saturated, sharp, commutative, unital semigroup.  Monoids are written additively.  A monoid $M$ induces a translation invariant partial order on its associated group $M^{\rm gp}$ such that $x \geq 0$ if and only if $x \in M$.  A valuative monoid is a monoid in which the associated partial order is a total order.

If $M$ is a monoid, $x \in M^{\rm gp}$, and $y \in M$, we write $x \prec y$ to mean that there are integers $m$ and $n$ such that $my \leq x \leq ny$.



A monoid homomorphism $f : M \to N$ is called \emph{infinitesimal} if the following condition holds:
\begin{enumerate}[label=$\bullet$]
\item If there are $x \in M^{\rm gp}$ and $y \in M$ such that $f(x) \in N$ and $x \prec y$ and $f(y) \prec f(x)$ then $x \in M$.
\end{enumerate}

If $P \xrightarrow f Q \xrightarrow g N$ are monoid homomorphisms, we say that $f$ is \emph{birhizomic} at $g$ if it is of finite presentation in the category of (integral, saturated, sharp) monoids and, for every factorization of $gf$ through monoid homomorphisms $P \xrightarrow h M \xrightarrow p N$ with $p$ infinitesimal, there is a unique monoid homomorphism $u : Q \to M$ with $pu = g$ and $uf = h$.

A family of morphisms $P \rightarrow Q_i$ is called a birhizomic cover if, for every valuative monoid $V$ and every homomorphism $v : P \to V$, there is a factorization of $v$ through some $Q_i \to V$ such that $P \to Q_i$ is birhizomic at $Q_i \to V$.

Let $P$ be a monoid.  Let $A_P$ be the category of monoid homomorphisms $P \to Q$.  The birhizomic covers form a Grothendieck topology on $A_P^{\rm op}$.



Let $P$ be a monoid.  For any $x$ in $P$, there is an initial monoid homomorphism $P \to P[-x]^\sharp$ that sends $x$ to $0$.  Let $B_P$ be the full subcategory of $A_P$ consisting of homomorphisms $P \to P[-x]^\sharp$, for $x \in P$.

We have given $A_P^{\rm op}$ the birhizomic topology and $B_P^{\rm op}$ the trivial topology in which all presheaves are sheaves.  The inclusion $B_P^{\rm op} \subset A_P^{\rm op}$ is left exact.  By left Kan extension and sheafification, it extends to a functor $\pi^{-1} : \operatorname{Sh}(B_P) \to \operatorname{Sh}(A_P)$ on the categories of sheaves.   Let $\pi_\ast$ be the right adjoint of $\pi^{-1}$.



\paragraph{Question}

\em
Is it the case that, for all presheaves of abelian groups $F$ on $B_P^{\rm op}$, the morphism $$F \to \mathrm R \pi_\ast \pi^{-1} F$$ is a quasi-isomorphism?
\end{document}
